\chapter{§2.7 二元一次方程组}
\label{sec:2.7}


\prerequisite{§2.6 一元一次方程}

\gradelevel{8 年级}


\section*{二元一次方程与方程组}


\begin{explore}


学校组织春游，包了大巴和中巴各若干辆。大巴每辆坐 $45$ 人，中巴每辆坐 $30$ 人，共包了 $5$ 辆车，恰好坐满 $180$ 人。大巴和中巴各多少辆？

设大巴 $x$ 辆，中巴 $y$ 辆。我们可以列出两个条件：

\begin{itemize}
  \item 车辆总数： $x + y = 5$
  \item 人数总数： $45x + 30y = 180$
\end{itemize}


一个方程不够确定两个未知数，但两个方程放在一起就能求解——这就是\textbf{方程组}。

\end{explore}

\begin{understand}


\textbf{二元一次方程}：含有两个未知数，且每个未知数的最高次数为 $1$ 的方程。如 $x + y = 5$ ， $2x - 3y = 7$ 。

\textbf{二元一次方程组}：由两个（或多个）二元一次方程组成的一组方程，要求同时满足。

\[
\begin{cases} x + y = 5 \\ 45x + 30y = 180 \end{cases}
\]


\textbf{方程组的解}：使方程组中每个方程都成立的一组未知数的值。

\textbf{为什么需要方程组？} 一个二元一次方程有无穷多组解（如 $x + y = 5$ 的解有 $(1,4), (2,3), \ldots$ ）。只有再加一个方程，才能把解"锁定"到唯一一组。

\end{understand}

\begin{workedexamples}


\textbf{例 1.} 检验 $x = 2, y = 3$ 是否为方程组 $\begin{cases} x + y = 5 \\ 45x + 30y = 180 \end{cases}$ 的解。

\textbf{解：}

\begin{itemize}
  \item 代入第一个方程： $2 + 3 = 5$ ✓
  \item 代入第二个方程： $45 \times 2 + 30 \times 3 = 90 + 90 = 180$ ✓
\end{itemize}


两个方程都成立，所以 $x = 2, y = 3$ 是方程组的解。

\textbf{例 2.} 检验 $x = 3, y = 2$ 是否为方程组 $\begin{cases} x + y = 5 \\ 45x + 30y = 180 \end{cases}$ 的解。

\textbf{解：}

\begin{itemize}
  \item 第一个方程： $3 + 2 = 5$ ✓
  \item 第二个方程： $45 \times 3 + 30 \times 2 = 135 + 60 = 195 \neq 180$ ✗
\end{itemize}


第二个方程不成立，所以这不是方程组的解。

\textbf{例 3.} 写出一个以 $x = 1, y = 2$ 为解的二元一次方程组。

\textbf{解：} 例如 $\begin{cases} x + y = 3 \\ 2x - y = 0 \end{cases}$ 。验证： $1 + 2 = 3$ ✓， $2 \times 1 - 2 = 0$ ✓。

\end{workedexamples}

\begin{keytakeaway}


\begin{itemize}
  \item 二元一次方程组由两个含两个未知数的一次方程组成。
  \item 方程组的解必须同时满足所有方程。
  \item 一般情况下，两个方程确定两个未知数（有唯一解）。
\end{itemize}


\end{keytakeaway}

\begin{practice}


\begin{enumerate}
  \item 判断 $\begin{cases} x - y = 1 \\ 2x + y = 8 \end{cases}$ 中 $x = 3, y = 2$ 是否为解。
  \item 方程 $2x + y = 10$ 中，若 $x = 3$ ，则 $y$ 等于多少？
  \item 写出一个以 $x = 2, y = -1$ 为解的二元一次方程组。
  \item 一个二元一次方程有多少组解？
  \item 为什么单独一个二元一次方程无法确定唯一解？
\end{enumerate}


\end{practice}

\section*{代入消元法}


\begin{explore}


解方程组的核心思想是\textbf{消元}——把两个未知数的问题变成一个未知数的问题。第一种消元方法叫"代入法"：从一个方程中把一个未知数用另一个表示，然后代入另一个方程。

\end{explore}

\begin{understand}


\textbf{代入消元法}的步骤：

\begin{enumerate}
  \item 从某个方程中，将一个未知数用另一个未知数表示。
  \item 将这个表达式代入另一个方程，消去一个未知数。
  \item 解一元一次方程，求出一个未知数。
  \item 将求得的值代回，求出另一个未知数。
\end{enumerate}


\end{understand}

\begin{workedexamples}


\textbf{例 1.} 用代入法解方程组 $\begin{cases} y = 2x - 1 \quad \cdots ① \\ 3x + 2y = 5 \quad \cdots ② \end{cases}$

\textbf{解：}

①已经将 $y$ 用 $x$ 表示。将 $y = 2x - 1$ 代入②：

\[
3x + 2(2x - 1) = 5
\]


\[
3x + 4x - 2 = 5
\]


\[
7x = 7
\]


\[
x = 1
\]


将 $x = 1$ 代入①： $y = 2(1) - 1 = 1$ 。

解为 $\begin{cases} x = 1 \\ y = 1 \end{cases}$ 。

\textbf{例 2.} 用代入法解方程组 $\begin{cases} x + 2y = 7 \quad \cdots ① \\ 3x - y = 7 \quad \cdots ② \end{cases}$

\textbf{解：}

从①中解出 $x$ ： $x = 7 - 2y$ ...... ③

将③代入②： $3(7 - 2y) - y = 7$ 。

\[
21 - 6y - y = 7
\]


\[
-7y = -14
\]


\[
y = 2
\]


将 $y = 2$ 代入③： $x = 7 - 2(2) = 3$ 。

解为 $\begin{cases} x = 3 \\ y = 2 \end{cases}$ 。

\textbf{例 3.} 用代入法解方程组 $\begin{cases} 2x + 3y = 12 \quad \cdots ① \\ x = y + 1 \quad \cdots ② \end{cases}$

\textbf{解：}

将②代入①： $2(y + 1) + 3y = 12$ 。

\[
2y + 2 + 3y = 12
\]


\[
5y = 10
\]


\[
y = 2
\]


$x = 2 + 1 = 3$ 。

解为 $\begin{cases} x = 3 \\ y = 2 \end{cases}$ 。

\end{workedexamples}

\begin{keytakeaway}


\begin{itemize}
  \item 代入法适合某个方程中某个未知数系数为 $1$ 或 $-1$ 的情况。
  \item 关键是"用一个表示另一个"，再代入另一个方程消元。
  \item 求出一个未知数后，代回求另一个。
\end{itemize}


\end{keytakeaway}

\begin{practice}


\begin{enumerate}
  \item 用代入法解 $\begin{cases} y = x + 3 \\ 2x + y = 9 \end{cases}$ 。
  \item 用代入法解 $\begin{cases} x - y = 4 \\ 3x + 2y = 7 \end{cases}$ 。
  \item 用代入法解 $\begin{cases} 2x + y = 5 \\ x + 3y = 10 \end{cases}$ 。
  \item 用代入法解 $\begin{cases} x = 3y \\ 2x - y = 10 \end{cases}$ 。
  \item 为什么代入法适合系数为 $1$ 的情况？
\end{enumerate}


\end{practice}

\section*{加减消元法}


\begin{explore}


对于 $\begin{cases} 3x + 2y = 16 \\ 3x - 2y = 8 \end{cases}$ ，两个方程中 $x$ 的系数都是 $3$ ， $y$ 的系数一正一负。如果把两个方程直接\textbf{相加}， $2y$ 和 $-2y$ 就抵消了！这就是加减消元法的思路。

\end{explore}

\begin{understand}


\textbf{加减消元法}的步骤：

\begin{enumerate}
  \item 必要时，将某个方程的两边乘以适当的数，使两个方程中某个未知数的系数\textbf{相等或相反}。
  \item 两个方程\textbf{相加}（系数相反时）或\textbf{相减}（系数相等时），消去一个未知数。
  \item 解一元一次方程，求出一个未知数。
  \item 代回求另一个未知数。
\end{enumerate}


\end{understand}

\begin{workedexamples}


\textbf{例 1.} 用加减法解方程组 $\begin{cases} 3x + 2y = 16 \quad \cdots ① \\ 3x - 2y = 8 \quad \cdots ② \end{cases}$

\textbf{解：}

① $+$ ②： $6x = 24$ ， $x = 4$ 。

将 $x = 4$ 代入①： $12 + 2y = 16$ ， $2y = 4$ ， $y = 2$ 。

解为 $\begin{cases} x = 4 \\ y = 2 \end{cases}$ 。

\textbf{例 2.} 用加减法解方程组 $\begin{cases} 2x + 3y = 13 \quad \cdots ① \\ 5x + 3y = 22 \quad \cdots ② \end{cases}$

\textbf{解：}

$y$ 的系数相同，② $-$ ①： $3x = 9$ ， $x = 3$ 。

代入①： $6 + 3y = 13$ ， $3y = 7$ ， $y = \frac{7}{3}$ 。

解为 $\begin{cases} x = 3 \\ y = \frac{7}{3} \end{cases}$ 。

\textbf{例 3.} 用加减法解方程组 $\begin{cases} 3x + 4y = 10 \quad \cdots ① \\ 2x + 3y = 7 \quad \cdots ② \end{cases}$

\textbf{解：}

为消去 $x$ ，① $\times 2$ ，② $\times 3$ ：

\[
\begin{cases} 6x + 8y = 20 \quad \cdots ③ \\ 6x + 9y = 21 \quad \cdots ④ \end{cases}
\]


④ $-$ ③： $y = 1$ 。

代入②： $2x + 3 = 7$ ， $2x = 4$ ， $x = 2$ 。

解为 $\begin{cases} x = 2 \\ y = 1 \end{cases}$ 。

\textbf{例 4.} 用加减法解方程组 $\begin{cases} 4x - 5y = -3 \quad \cdots ① \\ 3x + 2y = 12 \quad \cdots ② \end{cases}$

\textbf{解：}

为消去 $y$ ，① $\times 2$ ，② $\times 5$ ：

\[
\begin{cases} 8x - 10y = -6 \quad \cdots ③ \\ 15x + 10y = 60 \quad \cdots ④ \end{cases}
\]


③ $+$ ④： $23x = 54$ ， $x = \frac{54}{23}$ 。

代入②： $3 \times \frac{54}{23} + 2y = 12$ ， $\frac{162}{23} + 2y = 12$ ， $2y = \frac{276 - 162}{23} = \frac{114}{23}$ ， $y = \frac{57}{23}$ 。

解为 $\begin{cases} x = \frac{54}{23} \\ y = \frac{57}{23} \end{cases}$ 。

\end{workedexamples}

\begin{keytakeaway}


\begin{itemize}
  \item 加减法的核心：让某个未知数的系数相等或相反，然后加减消元。
  \item 系数相反时相加，系数相等时相减。
  \item 当两种方法都可用时，选择计算更简便的一种。
\end{itemize}


\end{keytakeaway}

\begin{practice}


\begin{enumerate}
  \item 用加减法解 $\begin{cases} x + y = 10 \\ x - y = 4 \end{cases}$ 。
  \item 用加减法解 $\begin{cases} 2x + y = 7 \\ 3x + y = 9 \end{cases}$ 。
  \item 用加减法解 $\begin{cases} 2x + 3y = 8 \\ 3x - 2y = 1 \end{cases}$ 。
  \item 用加减法解 $\begin{cases} 5x - 2y = 1 \\ 3x + 4y = 11 \end{cases}$ 。
  \item 什么时候用加减法比代入法更方便？
  \item 用加减法解 $\begin{cases} 4x + 3y = 11 \\ 2x - y = 1 \end{cases}$ 。
\end{enumerate}


\end{practice}

\section*{三元一次方程组（简单）}


\begin{explore}


商店里， $2$ 斤苹果、 $3$ 斤橘子、 $1$ 斤香蕉共 $28$ 元； $1$ 斤苹果、 $2$ 斤橘子、 $2$ 斤香蕉共 $23$ 元； $3$ 斤苹果、 $1$ 斤橘子、 $3$ 斤香蕉共 $32$ 元。三种水果各多少钱一斤？这需要三个方程来确定三个未知数。

\end{explore}

\begin{understand}


\textbf{三元一次方程组}：含有三个未知数的一次方程组（通常需要三个方程）。

解法：逐步消元，先用两两搭配的方式消去一个未知数，化为二元一次方程组，再继续消元。

\end{understand}

\begin{workedexamples}


\textbf{例 1.} 解方程组

\[
\begin{cases}
x + y + z = 6 \quad \cdots ① \\
2x + y - z = 1 \quad \cdots ② \\
x - y + 2z = 7 \quad \cdots ③
\end{cases}
\]


\textbf{解：}

① $+$ ②： $3x + 2y = 7$ ...... ④

① $\times 2 +$ ③： $2x + 2y + 2z + x - y + 2z = 12 + 7$ ，即 $3x + y + 4z = 19$ 。

换个思路。用 ① 和 ② 消去 $z$ ：① $+$ ②： $3x + 2y = 7$ ...... ④

用 ① 和 ③ 消去 $z$ ：① $\times 2 -$ ③： $(2x + 2y + 2z) - (x - y + 2z) = 12 - 7$ ，得 $x + 3y = 5$ ...... ⑤

由 ④ 和 ⑤： $\begin{cases} 3x + 2y = 7 \\ x + 3y = 5 \end{cases}$

⑤ $\times 3$ ： $3x + 9y = 15$ ...... ⑥

⑥ $-$ ④： $7y = 8$ ， $y = \frac{8}{7}$ 。

代入⑤： $x = 5 - 3 \times \frac{8}{7} = 5 - \frac{24}{7} = \frac{11}{7}$ 。

代入①： $z = 6 - \frac{11}{7} - \frac{8}{7} = 6 - \frac{19}{7} = \frac{23}{7}$ 。

解为 $x = \frac{11}{7}$ ， $y = \frac{8}{7}$ ， $z = \frac{23}{7}$ 。

\textbf{例 2.} 解方程组

\[
\begin{cases}
x + y = 5 \quad \cdots ① \\
y + z = 7 \quad \cdots ② \\
x + z = 6 \quad \cdots ③
\end{cases}
\]


\textbf{解：} ① $+$ ② $+$ ③： $2(x + y + z) = 18$ ，所以 $x + y + z = 9$ ...... ④

④ $-$ ②： $x = 2$ 。④ $-$ ③： $y = 3$ 。④ $-$ ①： $z = 4$ 。

解为 $x = 2$ ， $y = 3$ ， $z = 4$ 。

\end{workedexamples}

\begin{keytakeaway}


\begin{itemize}
  \item 三元方程组的解法是逐步消元：三元 → 二元 → 一元。
  \item 注意选择合适的方程对进行消元，使计算简便。
  \item 特殊结构的方程组（如每个方程恰含两个未知数）可用巧妙方法求解。
\end{itemize}


\end{keytakeaway}

\begin{practice}


\begin{enumerate}
  \item 解方程组 $\begin{cases} x + y + z = 10 \\ x - y = 2 \\ y - z = 1 \end{cases}$ 。
  \item 解方程组 $\begin{cases} x + y = 4 \\ y + z = 6 \\ x + z = 8 \end{cases}$ 。
  \item 解方程组 $\begin{cases} 2x + y - z = 3 \\ x - y + z = 5 \\ 3x + 2y - z = 6 \end{cases}$ 。
\end{enumerate}


\end{practice}

\section*{方程组的应用}


\begin{explore}


一元一次方程只有一个未知数，列方程时需要巧妙地将多个未知量用一个表达式关联。有了方程组，我们可以直接设多个未知数，更自然地表达问题中的数量关系。

\end{explore}

\begin{understand}


列方程组解应用题的步骤与列一元一次方程类似，但可以设多个未知数：

\begin{enumerate}
  \item 审题，确定未知量。
  \item 设两个或多个未知数。
  \item 找出与未知数个数相同的等量关系。
  \item 列方程组，求解。
  \item 检验作答。
\end{enumerate}


\end{understand}

\begin{workedexamples}


\textbf{例 1.} 一支钢笔和两支圆珠笔共 $13$ 元，两支钢笔和一支圆珠笔共 $14$ 元。求钢笔和圆珠笔各多少元？

\textbf{解：} 设钢笔 $x$ 元，圆珠笔 $y$ 元。

\[
\begin{cases} x + 2y = 13 \quad \cdots ① \\ 2x + y = 14 \quad \cdots ② \end{cases}
\]


② $\times 2$ ： $4x + 2y = 28$ ...... ③

③ $-$ ①： $3x = 15$ ， $x = 5$ 。

代入①： $5 + 2y = 13$ ， $y = 4$ 。

答：钢笔 $5$ 元，圆珠笔 $4$ 元。

\textbf{例 2.} 甲、乙两人骑车从 A 地到 B 地。甲的速度为每小时 $15$ 千米，乙的速度为每小时 $12$ 千米。甲比乙晚出发 $1$ 小时，但比乙早到 $0.5$ 小时。求 A、B 两地的距离。

\textbf{解：} 设 A、B 两地距离为 $s$ 千米，甲用时 $t$ 小时。

\begin{itemize}
  \item $s = 15t$ （甲的路程等式）
  \item 乙用时 $(t + 1 + 0.5)$ 小时： $s = 12(t + 1.5)$
\end{itemize}


\[
15t = 12(t + 1.5)
\]


\[
15t = 12t + 18
\]


\[
3t = 18, \quad t = 6
\]


\[
s = 15 \times 6 = 90
\]


答：A、B 两地距离 $90$ 千米。

\textbf{例 3.} 一个两位数，十位数字与个位数字之和为 $9$ ，如果将十位与个位数字交换，得到的新数比原数大 $27$ 。求原来的两位数。

\textbf{解：} 设十位数字为 $x$ ，个位数字为 $y$ 。原数为 $10x + y$ ，新数为 $10y + x$ 。

\[
\begin{cases} x + y = 9 \\ (10y + x) - (10x + y) = 27 \end{cases}
\]


第二个方程化简： $9y - 9x = 27$ ，即 $y - x = 3$ ...... ②

由 $x + y = 9$ 和 $y - x = 3$ ：

两式相加： $2y = 12$ ， $y = 6$ ， $x = 3$ 。

原来的两位数是 $36$ 。

验证： $63 - 36 = 27$ ✓。

\end{workedexamples}

\begin{keytakeaway}


\begin{itemize}
  \item 方程组让我们可以设多个未知数，使列方程更直接。
  \item 有几个未知数就需要找几个独立的等量关系。
  \item 解题后要检验答案是否满足题目的所有条件。
\end{itemize}


\end{keytakeaway}

\begin{practice}


\begin{enumerate}
  \item 一件上衣和两条裤子共 $250$ 元，两件上衣和一条裤子共 $290$ 元。上衣和裤子各多少元？
  \item 小明和小红共有 $50$ 本书，如果小明给小红 $5$ 本，两人的书一样多。两人原来各有多少本？
  \item 两个数之和为 $100$ ，其差为 $24$ 。这两个数分别是多少？
  \item 一个长方形的周长为 $28$ 厘米，长比宽多 $2$ 厘米，求长和宽。
  \item 鸡兔同笼：笼中有鸡和兔共 $20$ 只，脚共 $56$ 只。鸡和兔各多少只？
\end{enumerate}


\end{practice}



\begin{answer}


\textbf{二元一次方程与方程组 练一练}

\begin{enumerate}
  \item $3 - 2 = 1$ ✓， $2 \times 3 + 2 = 8$ ✓。是解。
  \item $2 \times 3 + y = 10$ ， $y = 4$ 。
  \item 例如 $\begin{cases} x + y = 1 \\ x - y = 3 \end{cases}$ 。
  \item 无穷多组解。
  \item 一个方程含两个未知数，任给一个未知数的值都能得到另一个，因此有无穷多组解。
\end{enumerate}


\textbf{代入消元法 练一练}

\begin{enumerate}
  \item 代入： $2x + (x + 3) = 9$ ， $3x = 6$ ， $x = 2$ ， $y = 5$ 。
  \item $x = y + 4$ 代入： $3(y + 4) + 2y = 7$ ， $5y = -5$ ， $y = -1$ ， $x = 3$ 。
  \item 由第一式： $y = 5 - 2x$ 代入第二式： $x + 3(5 - 2x) = 10$ ， $x + 15 - 6x = 10$ ， $-5x = -5$ ， $x = 1$ ， $y = 3$ 。
  \item 代入： $2(3y) - y = 10$ ， $5y = 10$ ， $y = 2$ ， $x = 6$ 。
  \item 系数为 $1$ 时可以直接解出一个未知数的表达式，避免分数运算。
\end{enumerate}


\textbf{加减消元法 练一练}

\begin{enumerate}
  \item 两式相加： $2x = 14$ ， $x = 7$ ， $y = 3$ 。
  \item ② $-$ ①： $x = 2$ ，代入①： $y = 3$ 。
  \item ① $\times 2$ ： $4x + 6y = 16$ ；② $\times 3$ ： $9x - 6y = 3$ 。相加： $13x = 19$ ， $x = \frac{19}{13}$ ； $y = \frac{8 - 2 \times \frac{19}{13}}{3} = \frac{8 - \frac{38}{13}}{3} = \frac{\frac{66}{13}}{3} = \frac{22}{13}$ 。
  \item ① $\times 2$ ： $10x - 4y = 2$ ，加②： $13x = 13$ ， $x = 1$ ，代入①： $y = 2$ 。
  \item 当两个方程中某个未知数的系数相同或相反时，加减法更方便。
  \item ② $\times 3$ ： $6x - 3y = 3$ ，加①： $10x = 14$ ， $x = \frac{7}{5}$ 。代入②： $y = 2 \times \frac{7}{5} - 1 = \frac{9}{5}$ 。
\end{enumerate}


\textbf{三元一次方程组 练一练}

\begin{enumerate}
  \item 由②： $x = y + 2$ ，由③： $z = y - 1$ 。代入①： $(y + 2) + y + (y - 1) = 10$ ， $3y + 1 = 10$ ， $y = 3$ ， $x = 5$ ， $z = 2$ 。
  \item 三式相加： $2(x + y + z) = 18$ ， $x + y + z = 9$ 。 $z = 9 - 4 = 5$ ， $x = 9 - 6 = 3$ ， $y = 9 - 8 = 1$ 。
  \item ① $+$ ②： $3x = 8$ ， $x = \frac{8}{3}$ 。① $-$ ③（不行，需要先消元）。①和③消去 $z$ ：① $-$ ③ 不可（系数相同），直接相减： $(2x + y - z) - (3x + 2y - z) = 3 - 6$ ， $-x - y = -3$ ，即 $x + y = 3$ ...... ④。由 ① $+$ ②： $3x = 8$ ， $x = \frac{8}{3}$ ， $y = 3 - \frac{8}{3} = \frac{1}{3}$ 。代入②： $\frac{8}{3} - \frac{1}{3} + z = 5$ ， $z = 5 - \frac{7}{3} = \frac{8}{3}$ 。解为 $x = \frac{8}{3}$ ， $y = \frac{1}{3}$ ， $z = \frac{8}{3}$ 。
\end{enumerate}


\textbf{方程组的应用 练一练}

\begin{enumerate}
  \item 设上衣 $x$ 元，裤子 $y$ 元。 $\begin{cases} x + 2y = 250 \\ 2x + y = 290 \end{cases}$ 。两式相加： $3x + 3y = 540$ ， $x + y = 180$ 。由 $x + 2y = 250$ 减去 $x + y = 180$ ： $y = 70$ ， $x = 110$ 。上衣 $110$ 元，裤子 $70$ 元。
  \item 设小明 $x$ 本，小红 $y$ 本。 $\begin{cases} x + y = 50 \\ x - 5 = y + 5 \end{cases}$ 。由②： $x - y = 10$ 。解得 $x = 30$ ， $y = 20$ 。
  \item $\begin{cases} x + y = 100 \\ x - y = 24 \end{cases}$ ， $x = 62$ ， $y = 38$ 。
  \item $\begin{cases} 2(l + w) = 28 \\ l - w = 2 \end{cases}$ ，即 $\begin{cases} l + w = 14 \\ l - w = 2 \end{cases}$ 。 $l = 8$ ， $w = 6$ 。长 $8$ 厘米，宽 $6$ 厘米。
  \item 设鸡 $x$ 只，兔 $y$ 只。 $\begin{cases} x + y = 20 \\ 2x + 4y = 56 \end{cases}$ 。由②： $x + 2y = 28$ 。减①： $y = 8$ ， $x = 12$ 。鸡 $12$ 只，兔 $8$ 只。
\end{enumerate}


\end{answer}
